\section{The curvature-kernel field and a reduced flat alteration}\label{har:sec:alteration}

This construction converts ordinary cohomology on $X$ itself; in the descent application, $X$ will be a finite \'etale cover. It does not require \(\operatorname{Pic}^0=0\).
Its alteration is genuinely generically finite. We will explicitly
descend the resulting metrized base line to birational models of the
original function field.

\begin{theorem}[Ordinary conversion through a curvature-null field]\label{thm:alteration-conversion}
Let $X$ be smooth connected projective, with smoothly semipositive $L=-K_X$. Suppose that for one fixed $q\ge0$, $H^q(X,mL)\ne0$ for unbounded positive integers $m$. Then $H^0(X,rL)\ne0$ for some $r>0$.
\end{theorem}

Smooth Lefschetz and Corollary~\ref{inj:forms-input} already imply this
conclusion. We give a different geometric proof through three
constructions: the field constant in curvature-null directions, a line
of multiplication maps on a reduced flat alteration, and descent of
that metrized line and its injection to the original function field.
The alteration makes the fiberwise maximum positive and continuous by
using reduced fibers. Returning to the original function field then
allows the birational Euler argument to apply.

\subsection{The field constant in curvature-null directions}\label{har:subsec:kernelfield}

Let \(X\) be smooth connected projective with smoothly semipositive
\(L=-K_X\), and suppose \(H^q(X,mL)\ne0\) for arbitrarily large
positive \(m\), with fixed \(q\). Put \(n=\dim X\), let \(\theta\)
be the smooth semipositive curvature of \(L\), and let \(r\) be its
maximal rank. If \(q=0\), the conclusion is already given. Otherwise
start from the assumed classes and apply the harmonic-star construction
in the proof of Lemma~\ref{har:lem:forms}; it gives nonzero sections
of \(\Omega_X^{n-q}\otimes L^{m+1}\). No invariant-index assumption
is involved. If \(q=n\), these are anticanonical sections; if
\(r=n\), the positive curvature energy on an open set excludes
\(q>0\). In the remaining cases the same determinant extraction gives
a line \(M\) mapping nontrivially to a positive exterior power of
\(\Omega_X^{n-q}\), and effective divisors
\[
 D_i\in|M+a_iL|,\qquad a_i\to\infty,
 \qquad D_i[\theta]^r[\omega]^{n-r-1}=0.
\]
The intersection equality is Lemma~\ref{har:lem:null-divisors}.
If \(r=0\), all \(D_i\) vanish, and two distinct \(a_i\) show that
a positive power of \(L\) is trivial. Assume \(1\le r<n\).

On the ordinary open \(U\) of maximal curvature rank, define
\[
 \mathcal K=\{g\in\mathbb C(X): dg(\ker\theta)=0
       \text{ on }U\text{ off a proper algebraic subset}\}.
\]
It is a field. It is relatively algebraically closed: differentiate a
minimal polynomial for an element algebraic over \(\mathcal K\) in a
null direction; its separable derivative is generically nonzero, so
the element has zero derivative there as well. The intermediate-field
argument from Section~\ref{har:subsec:small} proves finite generation.

If a rational function has null pole support, it belongs to
\(\mathcal K\). Its level divisors have the same zero intersection
as the poles, hence each of their components is null. At a regular
level point in \(U\), the tangent hyperplane contains \(\ker\theta\),
by the same rank-quotient argument as before. This proves the assertion
where the differential is nonzero and then everywhere required.

Choose a smooth projective model \(Y\) of \(\mathcal K\) and resolve
the map by \(\mu:X_1\to X\), \(f:X_1\to Y\). Put \(A=\mu^*L\).
The generic fiber is geometrically integral, because the extension
is regular in characteristic zero. On a nonempty ordinary open,
\(\ker\mu^*\theta\subset\ker df\), and \(\dim Y\le r\).
In fact for every smooth \((1,1)\)-form \(\eta\) on \(Y\),
\begin{equation}\label{har:eq:globalwedge}
 (\mu^*\theta)^r\wedge f^*\eta=0.
\end{equation}
On the rank-\(r\) open this follows because every horizontal
covector lies in the positive-curvature span; elsewhere the rank is
smaller, and continuity extends the identity across the algebraic
exceptional locus.

Any integral relation \(\sum b_i=\sum b_i a_i=0\) makes
\(\sum b_iD_i\) principal with null pole support. Its function is in
\(\mathcal K\), so the horizontal parts of \(\mu^*D_i\) satisfy the
same relation. Three-term interpolation proves that these horizontal
divisors are affine in \(a_i\). Their slope is effective by unbounded
nonnegative coefficients. Thus for two indices with \(b=a_2-a_1>0\)
the quotient \(s\) is a rational section of \(bA\) with effective
horizontal divisor, regular over a dense base open. If \(Y\) is a
point it already descends to a section on \(X\).

Let \(E=K_{X_1/X}=K_{X_1}+A\). For every \(t\ge0\),
\begin{equation}\label{har:eq:ordinaryrank}
 \operatorname{rank}f_*\mathcal O_{X_1}(E+tA)\le1,
 \quad\text{with equality for }t=kb,\ k\ge0.
\end{equation}
For the upper bound, twist two hypothetical independent generic
sections by a high ample base line to make them global. Their zero
divisors push to effective divisors on \(X\) with zero intersection
against \([\theta]^r[\omega]^{n-r-1}\), using
\eqref{har:eq:globalwedge}, exceptionality of \(E\), and
\(\theta^{r+1}=0\). Their quotient has null poles, hence lies in
\(\mathcal K\), contrary to independence. The lower bound is supplied
by \(s_Es^k\).

Over a smooth dense open with geometrically integral fibers, the
function \(-\log\max_{X_{1,y}}|s|^2\) is plurisubharmonic. Indeed the
rank-one adjoint section \((s_E/f^*\eta)s^k\) has smooth coefficient
\((kb+1)A\). The moment calculation
\eqref{har:eq:moments} applies, with full-support density and locally
uniform power-mean limit. The next construction extends this good-locus
metric using reduced flat fibers rather than divisorial minima.

\subsection{A Hom line on a reduced flat family}\label{har:subsec:homline}

By weak semistable reduction \cite[Theorem~0.3]{AK2000},
there are a smooth projective alteration \(Y_a\to Y\) and a
modification \(Z_a\) of the main component of \(X_1\times_Y Y_a\)
such that \(p_a:Z_a\to Y_a\) is equidimensional, toroidal, and has
reduced fibers. Toroidal varieties are Cohen--Macaulay, so the
equidimensional flatness criterion over the smooth base makes
\(p_a\) flat. Write \(q_a:Z_a\to X_1\). The source \(Z_a\) need not
be smooth.

Take a sufficiently high power \(\mathcal P\) of a relatively very
ample line. Relative Serre vanishing and base change give locally
free bundles
\[
 \mathcal H=p_{a*}\mathcal P,\qquad
 \mathcal H'=p_{a*}(\mathcal P\otimes q_a^*A^b),
\]
commuting with base change, with surjective evaluation for
\(\mathcal P\). On the generic fiber, multiplication by \(s\) gives
a nonzero map \(\mathcal H_\eta\to\mathcal H'_\eta\). Resolve its
rational point in the projective bundle of lines in
\(\operatorname{Hom}(\mathcal H,\mathcal H')\) on a smooth projective
birational base model \(Y_b\to Y_a\), and call the pulled-back
tautological subline \(B_b\). Put \(Z_b=Z_a\times_{Y_a}Y_b\), with
maps \(p_b\) and \(q_b\). Flatness and reduced geometric fibers
survive this base change. The total space is integral: flatness over
an integral base embeds its local rings in their generic localizations,
and its generic fiber is integral.

There is a regular line map
\begin{equation}\label{har:eq:homline}
 p_b^*B_b\longrightarrow q_b^*A^b
\end{equation}
nonzero on each fiber. To construct it, compose the tautological Hom
map with evaluation into \(\mathcal P\otimes q_b^*A^b\). It
annihilates the kernel of evaluation onto \(\mathcal P\), because
it is multiplication on the generic fiber and the total space is
integral. Factor through \(\mathcal P\) and cancel that invertible
line. If the result vanished as a map on a fiber, the induced map
of fiber section spaces would vanish, contrary to the everywhere
subbundle \(B_b\subset\operatorname{Hom}(\mathcal H,\mathcal H')\)
and base change. Reducedness now implies that it has a nonzero
value at some point of every fiber; without reducedness this last
implication would fail.

Give \(B_b\) the maximum metric through \eqref{har:eq:homline} and
the pulled-back smooth metric of \(A^b\). The maximum is positive
and continuous. Properness gives upper semicontinuity. Smooth points
are dense in each reduced fiber; a value arbitrarily close to the
maximum may be chosen at such a point. Flatness and fiber smoothness
make the morphism smooth there, giving liftings of nearby base points
and lower semicontinuity. On a dense open the map is the pullback of
\(s\) times a rational base factor, by its defining generic Hom line.
Its fibers map onto the original smooth integral fibers, so the
maximum weight equals the pullback of the good-locus weight just
proved, plus a pluriharmonic frame term. Continuity and removable
singularities extend plurisubharmonicity everywhere on \(Y_b\).

\subsection{Descending the alteration and its section}\label{har:subsec:alterationnorm}

Proposition~\ref{har:prop:euler} requires a source birational to \(X\),
whereas the construction above used an alteration. We first descend
a power of the metrized line by a finite group quotient. We then
descend its injection by a field norm. These two operations will
produce the required data on birational models of the original fields.

Choose a finite Galois extension of \(\mathbb C(Y)\) containing
\(\mathbb C(Y_b)\), of degree \(\ell\) and group \(G\). Starting
with the normalization of \(Y\) in this extension, normalize the simultaneous graph of
all conjugate rational maps to \(Y_b\). This gives a normal projective
\(W\), a regular \(G\)-action, and a map \(\rho:W\to Y_b\).
The quotient \(W_0=W/G\) is normal projective birational to \(Y\),
and \(W\to W_0\) is finite. Let \(B_W=\rho^*B_b\). The product
\[
 C_W=\bigotimes_{g\in G}g^*B_W
\]
has its permutation linearization and invariant product metric.
The power \(C_W^\ell\) has trivial stabilizer actions on fibers,
so descends to a line \(C_0\) on \(W_0\). One can check descent
locally on affine quotient neighborhoods: prescribe an equivariant
nonzero value on a finite orbit, lift it to a section, average, and
shrink the quotient neighborhood until it is a frame. Its invariant
transition functions descend. The invariant metric descends
continuously by the quotient topology.

Resolve \(\widehat Y\to W_0\), and let \(B_0\) be the pullback of
\(C_0\). Its metric is plurisubharmonic on the dense finite-étale
locus, and by continuity everywhere. Choose a smooth projective
birational model \(\widehat X\to X_1\) mapping to \(\widehat Y\),
with composite maps $\widehat\mu:\widehat X\to X$ and $\widehat f:\widehat X\to\widehat Y$, and put \(\widehat A=\widehat\mu^*L\). Normalize the main component of
\(\widehat X\times_{W_0}W\), calling it \(R\). It is finite over
\(\widehat X\) with the same Galois group: geometric integrality of
the generic fiber, proved above, makes the function field extension
linearly disjoint from the finite Galois base extension.
The finite maps fit into the commutative square
\[
\begin{array}{ccc}
 R & \longrightarrow & W\\
 \mathllap{\scriptstyle\mathrm{finite}\,}\big\downarrow &&
       \big\downarrow\mathrlap{\,\scriptstyle\mathrm{finite}}\\
 \widehat X & \longrightarrow & W_0 .
\end{array}
\]
Here the lower map factors through \(\widehat f:\widehat X\to
\widehat Y\), and the upper-right space maps to \(Y_b\) by \(\rho\).

On \(R\), the difference of the pullbacks \(b\widehat A-B_W\) has
a nonzero section. Indeed resolve the rational map from \(R\) to
\(Z_b\) over \(X_1\) and \(Y_b\). The resulting common graph is a
modification of \(R\), and \eqref{har:eq:homline} gives a regular
section there. Both line bundles already come from \(R\); normality
therefore descends this section to \(R\). Multiply its \(G\)-conjugates and
raise to the \(\ell\)-th power. This yields a section of the
pullback of \(b\ell^2\widehat A-\widehat f^*B_0\). A further field
norm for the finite surjection onto the normal target
\cite[Lemma~31.18.7]{StacksNormNormal}, of degree \(\ell\), relative to a rational line frame, gives
a regular nonzero section downstairs of
\begin{equation}\label{har:eq:descendedalteration}
 b\ell^3\widehat A-\widehat f^*(\ell B_0).
\end{equation}
Regularity can be checked at every prime: all orders upstairs are
nonnegative, hence so are the norm orders. The powers are not meant
to be minimal. Put \(\widehat b=b\ell^3\) and
\(\widehat Q=\ell B_0\). We have returned to birational models of
the original source and base fields, with a continuously semipositive
\(\widehat Q\) and an injection
\(\widehat f^*\widehat Q\to\widehat b\widehat A\).
The original kernel condition persists on a nonempty open of these
birational models. Theorem~\ref{har:thm:kernel} and
Proposition~\ref{har:prop:euler} now give a section of a positive
power of \(L\) on \(X\).
